
\section{Un motor impone campos, no una sola temperatura}
La temperatura del gas, la del metal y la de un sensor no son intercambiables. La refrigeración interna y superficial, el espesor, la conductividad y los recubrimientos crean gradientes. Un aumento de temperatura de gas no permite sumar el mismo incremento a todos los componentes. El cálculo térmico también necesita condiciones de contorno, transitorios y radiación cuando sea relevante.

Los gradientes generan deformaciones diferenciales; las restricciones convierten parte de ellas en tensión. En un arranque o parada pueden coexistir tensiones mecánicas y térmicas de signos y distribuciones cambiantes. La compatibilidad entre disco y álabe, holguras con carcasa y evolución del recubrimiento relaciona materiales y diseño térmico. Optimizar una propiedad aislada puede trasladar el problema a otra interfaz.
\section{Microestructura orientada, refrigeración y recubrimiento}
La solidificación direccional reduce ciertas fronteras transversales a la solicitación; un monocristal evita bordes de grano de alto ángulo dentro de la idealización del producto, pero conserva anisotropía cristalina y otros defectos posibles. Orientación, defectos de colada, tratamiento y acabado necesitan control. No equivale a un material sin mecanismos de fallo.

Un sistema de barrera térmica incluye funciones del recubrimiento cerámico, capa de unión y óxido que evoluciona durante exposición, según arquitectura. Diferencias de dilatación, oxidación, erosión y ciclos pueden iniciar pérdida de adherencia. Un recubrimiento desprendido modifica el campo térmico del sustrato; medir solo su espesor medio no describe todas las condiciones de protección. El ejemplo térmico P35 enseña resistencias en serie, sin reproducir una geometría refrigerada real.
\section{Un mapa de daño por componente}
En fan y compresor interesan impacto, erosión, fatiga y degradación aerodinámica; el material y la etapa condicionan la relevancia de cada uno. En cámaras y zonas calientes aparecen ciclos térmicos, oxidación, corrosión caliente, fisuración y degradación de recubrimientos. En discos se estudian fatiga y defectos en un campo centrífugo y térmico. En cojinetes, contacto, lubricación y contaminación exigen otro enfoque.

Este mapa se utiliza para formular preguntas de inspección: dónde, qué mecanismo, qué evidencia y qué método lo puede observar. No convierte una lista de daños posibles en diagnóstico. Dos mecanismos pueden producir señales parecidas; por ejemplo, una pérdida de rendimiento puede deberse a degradación de perfiles, holguras o condiciones de medida. El capítulo~\ref{ch:22} conecta este razonamiento con seguimiento en servicio.
\repaso{¿Por qué una cerámica puede trabajar en una zona caliente y aun así requerir protección ambiental?}{La estabilidad térmica no garantiza resistencia frente a todas las especies del gas, vapor, depósitos, erosión o ciclos. La selección considera química del ambiente y sistema completo de protección, además de la temperatura.}
